\PassOptionsToPackage{table,xcdraw}{xcolor}
\documentclass[aspectratio=169]{beamer}

\usetheme{Madrid}
\usecolortheme{default}
\definecolor{ucbblue}{RGB}{0, 51, 102}

\setbeamercolor{structure}{fg=ucbblue}
\setbeamercolor*{palette primary}{fg=white,bg=ucbblue}
\setbeamercolor*{palette secondary}{fg=white,bg=ucbblue!80}
\setbeamercolor*{palette tertiary}{fg=white,bg=ucbblue!90}
\setbeamercolor{title}{fg=white, bg=ucbblue}
\setbeamercolor{frametitle}{fg=white, bg=ucbblue}
\setbeamercolor{item}{fg=ucbblue}

\usepackage[T1]{fontenc}
\usepackage[utf8]{inputenc}
\usepackage{tgpagella}
\usefonttheme{serif}
\usepackage[spanish, es-noshorthands]{babel}
\usepackage{amsmath, amssymb, bm}
\usepackage{graphicx}
\usepackage[most]{tcolorbox}
\usepackage{booktabs}
\usepackage{tabularx}
\usepackage{tikz}
\usepackage[american]{circuitikz}
\usetikzlibrary{shapes.geometric, arrows.meta, positioning, calc}

\title[Par Diferencial DC]{Consolidación Transversal y Análisis DC del Par Diferencial BJT}
\subtitle{IMT-221 Circuitos Electrónicos III — Semana 7, Sesión 2}
\author[B. Quiroga Turdera]{M.Sc. Ing. Bernardo Quiroga Turdera}
\institute[UCB Tarija]{Universidad Católica Boliviana ``San Pablo'' — Sede Tarija}
\date{Semana 7}

\begin{document}

\begin{frame}
    \titlepage
\end{frame}

\begin{frame}{Objetivos de la Sesión}
    \textbf{Preguntas guía de hoy:}
    \begin{enumerate}
        \item ¿Podemos reconstruir el modelo de cualquier topología (CE, CC, CB) usando un solo método sin depender de fórmulas memorizadas?[cite: 15]
        \item ¿Qué cambia en el circuito cuando dos ramas de transistores ya no son independientes, sino que comparten una restricción de corriente (Par Diferencial)?[cite: 15]
    \end{enumerate}
    
    \vspace{0.4cm}
    \begin{tcolorbox}[colback=ucbblue!5, colframe=ucbblue, title=Transición Conceptual]
        Pasaremos del análisis de \textbf{un solo transistor} (donde la topología dicta la función) a una \textbf{etapa de detección acoplada de dos transistores} (el Par Diferencial)[cite: 15].
    \end{tcolorbox}
\end{frame}

\begin{frame}{Revisión Transversal: Un Método, Tres Topologías}
    No existen tres sets de fórmulas desconectadas. Existe \textbf{un solo método incremental}[cite: 15]:
    
    \vspace{0.4cm}
    \begin{center}
    \begin{tikzpicture}[auto, thick]
        % Uso de "positioning" para evitar solapamiento entre cajas, midiendo desde los bordes.
        \node[draw, fill=blue!10, rounded corners, align=center, text width=2cm] (q) {1. Punto Q \\ $I_{CQ}$};
        \node[draw, fill=blue!10, rounded corners, align=center, text width=2.4cm, right=0.4cm of q] (param) {2. Parámetros \\ $g_m, r_e, r_\pi$};
        \node[draw, fill=blue!10, rounded corners, align=center, text width=2.5cm, right=0.4cm of param] (ac) {3. Eq. AC \\ (Tierra AC)};
        \node[draw, fill=blue!10, rounded corners, align=center, text width=2.8cm, right=0.4cm of ac] (eq) {4. KVL / KCL \\ $A_v, R_{in}, R_{out}$};
        
        \draw[->] (q) -- (param);
        \draw[->] (param) -- (ac);
        \draw[->] (ac) -- (eq);
    \end{tikzpicture}
    \end{center}
    
    \vspace{0.5cm}
    \textbf{¿Qué es invariante?} La física del transistor incremental ($i_c = g_m v_\pi$)[cite: 15]. \\
    \textbf{¿Qué cambia?} El terminal común (Tierra AC), por dónde entra la señal ($v_i$) y por dónde sale ($v_o$)[cite: 15].
\end{frame}

\begin{frame}{Estructura del Par Diferencial Canónico}
    \begin{columns}
        \begin{column}{0.4\textwidth}
            \centering
            % Escala ligeramente ajustada para asegurar que no se solape con los bordes
            \begin{circuitikz}[scale=0.75, transform shape, thick]
                \draw (0,0) node[npn](Q1){}; \node[left=0.2cm] at (Q1.B) {$v_1$};
                \draw (3,0) node[npn, xscale=-1](Q2){}; \node[right=0.2cm] at (Q2.B) {$v_2$};
                \draw (Q1.C) to[R, l={$R_C$}, i={$I_{C1}$}] (0,2.5) -- (3,2.5) to[R, l_={$R_C$}, i_={$I_{C2}$}] (Q2.C);
                \draw (1.5,2.5) -- (1.5,3) node[above]{$+V_{CC}$};
                \draw (Q1.E) -- (0,-1.5) -- (3,-1.5) -- (Q2.E);
                \draw (1.5,-1.5) node[circle, fill, inner sep=1.5pt]{} node[right]{Nodo $E$} to[I, l={$I_{TAIL}$}] (1.5,-3.5) node[below]{$-V_{EE}$};
                \draw (Q1.B) to[short, -o] (-1,0);
                \draw (Q2.B) to[short, -o] (4,0);
                \draw (Q1.C) to[short, *-o] (-1,1) node[left]{$v_{C1}$};
                \draw (Q2.C) to[short, *-o] (4,1) node[right]{$v_{C2}$};
            \end{circuitikz}
        \end{column}
        \begin{column}{0.55\textwidth}
            \textbf{Anatomía de la etapa acoplada[cite: 15]:}
            \begin{itemize}
                \item \textbf{Dos ramas} ($Q_1$ y $Q_2$) idealmente idénticas (emparejadas).
                \item Dos nodos de entrada ($v_1, v_2$).
                \item Resistores de colector simétricos ($R_C$).
                \item \textbf{Restricción Fundamental:} Emisores unidos a una única fuente de corriente (Cola o \textit{Tail}) $I_{TAIL}$[cite: 15].
            \end{itemize}
            
            \vspace{0.2cm}
            \begin{alertblock}{No son dos CE independientes}
                Por KCL en el nodo E: $I_{E1} + I_{E2} = I_{TAIL}$[cite: 15]. ¡Ambas ramas deben \textbf{compartir} esta corriente constante!
            \end{alertblock}
        \end{column}
    \end{columns}
\end{frame}

\begin{frame}{Análisis DC: Estado Balanceado ($v_1 = v_2$)}
    Si aplicamos exactamente la misma tensión DC en ambas bases ($v_1 = v_2 = V_{CM}$) y los transistores son idénticos[cite: 15]:
    
    \vspace{0.2cm}
    \textbf{1. Simetría de Corrientes[cite: 15]:}
    Como $v_{BE1} = v_1 - V_E$ y $v_{BE2} = v_2 - V_E$, entonces $v_{BE1} = v_{BE2}$. Por lo tanto:
    \begin{equation*}
        I_{E1} = I_{E2} \quad \text{y como} \quad I_{E1} + I_{E2} = I_{TAIL} \implies I_{E1} = I_{E2} = \frac{I_{TAIL}}{2} \quad \text{[cite: 15]}
    \end{equation*}
    
    \textbf{2. Corrientes de Colector (Punto Q)[cite: 15]:}
    Asumiendo $\alpha \approx 1$ (despreciando corriente de base para este análisis simplificado):
    \begin{equation*}
        I_{C1} \approx I_{C2} \approx \frac{I_{TAIL}}{2} \quad \text{[cite: 15]}
    \end{equation*}
    
    \textbf{3. Tensiones de Colector (Salida DC)[cite: 15]:}
    Por KVL desde la fuente $V_{CC}$ hacia abajo:
    \begin{equation*}
        V_{C1} = V_{CC} - I_{C1}R_C \quad \text{y} \quad V_{C2} = V_{CC} - I_{C2}R_C \implies \mathbf{V_{C1} = V_{C2}} \quad \text{[cite: 15]}
    \end{equation*}
\end{frame}

\begin{frame}{Direccionamiento de Corriente (Current Steering)}
    ¿Qué ocurre si introducimos una diferencia $v_d = v_1 - v_2$?[cite: 15]
    
    \vspace{0.1cm}
    % Reducción de espaciado y tamaño para evitar Overfull \vbox
    \begin{table}
        \centering
        \small
        \renewcommand{\arraystretch}{1.2}
        \begin{tabularx}{0.95\textwidth}{|>{\centering\arraybackslash}p{2.5cm}|>{\centering\arraybackslash}X|>{\centering\arraybackslash}X|>{\centering\arraybackslash}X|>{\centering\arraybackslash}X|}
        \hline
        \rowcolor{ucbblue!20}
        \textbf{Condición} & $\mathbf{I_{C1}}$ & $\mathbf{I_{C2}}$ & $\mathbf{V_{C1}}$ (Salida 1) & $\mathbf{V_{C2}}$ (Salida 2) \\ \hline
        $\mathbf{v_1 = v_2}$ ($v_d = 0$) & $I_{TAIL}/2$ & $I_{TAIL}/2$ & Balanceado & Balanceado \\ \hline
        $\mathbf{v_1 > v_2}$ ($v_d > 0$) & $\uparrow$ (Sube) & $\downarrow$ (Baja) & $\downarrow$ (Cae) & $\uparrow$ (Sube) \\ \hline
        $\mathbf{v_1 < v_2}$ ($v_d < 0$) & $\downarrow$ (Baja) & $\uparrow$ (Sube) & $\uparrow$ (Sube) & $\downarrow$ (Cae) \\ \hline
        \end{tabularx}
    \end{table}
    
    \vspace{0.1cm}
    \begin{alertblock}{Física del Direccionamiento}
        La corriente $I_{TAIL}$ es un recurso constante. Un aumento en $v_1$ hace que $Q_1$ conduzca más, lo que \textbf{obliga} a $Q_2$ a conducir menos para cumplir $I_{E1} + I_{E2} = I_{TAIL}$[cite: 15]. Las tensiones de colector se mueven en \textbf{direcciones opuestas}[cite: 15].
    \end{alertblock}
\end{frame}

\begin{frame}{La Realidad del Diseño: Sensibilidad y "Mismatch"}
    ¿Por qué el balance perfecto $V_{C1} = V_{C2}$ casi nunca ocurre en hardware real aunque $v_1 = v_2$?[cite: 15]
    
    \vspace{0.3cm}
    \begin{itemize}
        \item \textbf{Mismatch de transistores:} $Q_1$ y $Q_2$ tienen leves diferencias de fabricación en el voltaje $V_{BE}$ y la ganancia $\beta$[cite: 15].
        \item \textbf{Mismatch de Resistores:} $R_{C1}$ y $R_{C2}$ tienen tolerancias reales (ej. 1\% o 5\%)[cite: 15].
        \item \textbf{Deriva Térmica:} Diferencias sutiles de temperatura alteran la distribución de corriente[cite: 15].
    \end{itemize}
    
    \vspace{0.3cm}
    Todo esto genera un \textbf{Voltaje de Offset} en la salida. Para aplicaciones de precisión (como medir una caída minúscula en un Shunt de corriente para el \textbf{Hito 2}), requeriremos transistores acoplados térmicamente y análisis formal de CMRR[cite: 15].
\end{frame}

\begin{frame}{Planificación de Laboratorio (Hito 2)}
    El \textbf{Hito 2} requerirá implementar un Par Diferencial discreto para detección de corriente mediante shunt[cite: 15].
    
    \vspace{0.3cm}
    \textbf{Aviso Operacional:}
    \begin{itemize}
        \item \textbf{Fecha objetivo provisional:} Semana 8, Sesión 2[cite: 15].
        \item \textbf{Lista de Materiales (BOM):} NO compren componentes aún. La topología exacta, las tensiones de alimentación y la especificación del transistor/shunt están pendientes de los documentos autoritativos del curso[cite: 15].
        \item El éxito en el laboratorio requerirá validar teóricamente el CMRR $> 60\,\text{dB}$, concepto matemático que exploraremos la próxima sesión[cite: 15].
    \end{itemize}
\end{frame}

\end{document}
